% GENERATED FILE. Edit canonical lessons, then run npm run generate:books.
% canonical-source-hash: fnv1a64:52b932b836d25c2a
% canonical-lessons: BN-C53-pora, BN-C53-chhara, BN-C53-phla, BN-C53-ora, BN-C53-shekhano

\chapter{To wear, To leave, To throw, To fly, To teach}
\label{ch:bn-to-wear-to-leave-to-throw-to-fly-to-teach}
\hlchaptermodality{53}

\begin{chapteropening}
\textbf{By the end of this chapter:} \emph{I can say to wear, to leave, to let go, to throw, to drop, to fly and to teach.}

Complete the last lesson of chapter 53: I can say to wear, to leave, to let go, to throw, to drop, to fly and to teach.
\end{chapteropening}

\section[\texorpdfstring{\bn{পরা} (pôrā)}{পরা (pôrā)}]{\bn{পরা} (pôrā) — to wear}

\label{lesson:BN-C53-pora}



\begin{quote}
Before the new one: say the Bengali for to look for, then the Bengali for to lose.
\end{quote}

\subsection*{You'll want to know: \bn{পরা}}
\textbf{\bn{পরা}} — \emph{pôrā} — \textquotedblleft{}to wear\textquotedblright{}.

\begin{cousinweb}
\textbf{\bn{পরা}} (\emph{pôrā}) is a tadbhava, a word inherited from Sanskrit and reshaped by sound change. It comes from \emph{paridadhāti}, \textquotedblleft{}sets round, puts on clothes\textquotedblright{} (CDIAL 7835). CDIAL notes that Chatterji preferred the passive \emph{paridhīyate}.
\end{cousinweb}

\begin{grammarlens}[title={What you've built}]
One more word for this chapter.
\end{grammarlens}

\subsection*{Your turn}
\begin{itemize}
\raggedright
  \item \emph{Say it:} \emph{pôrā}
  \item \emph{Say it:} \emph{pôrā}, once more
  \item \emph{Say it:} say \emph{pôrā}
  \item \emph{Recall:} say the Bengali for to look for, then the Bengali for to lose, then say \emph{pôrā} again
\end{itemize}

\begin{tcolorbox}[breakable,colback=teal!4,colframe=teal!35!black,title={Before you move on}]
What does \bn{পরা} mean? (\textquotedblleft{}To wear\textquotedblright{}.) Say it once more.
\end{tcolorbox}
\section[\texorpdfstring{\bn{ছাড়া} (chhāṛā)}{ছাড়া (chhāṛā)}]{\bn{ছাড়া} (chhāṛā) — to leave, to let go}

\label{lesson:BN-C53-chhara}



\begin{quote}
Before the new one: say the Bengali for to lose, then the Bengali for to wear.
\end{quote}

\subsection*{You'll want to know: \bn{ছাড়া}}
\textbf{\bn{ছাড়া}} — \emph{chhāṛā} — \textquotedblleft{}to leave, to let go\textquotedblright{}.

\begin{cousinweb}
\textbf{\bn{ছাড়া}} (\emph{chhāṛā}) is a tadbhava, a word inherited from Sanskrit and reshaped by sound change. It comes from \emph{chardayati}, \textquotedblleft{}causes to flow over, vomits\textquotedblright{} (CDIAL 4998). CDIAL gives its basic sense as \textquotedblleft{}eject\textquotedblright{}.
\end{cousinweb}

\begin{grammarlens}[title={What you've built}]
One more word for this chapter.
\end{grammarlens}

\subsection*{Your turn}
\begin{itemize}
\raggedright
  \item \emph{Say it:} \emph{chhāṛā}
  \item \emph{Say it:} \emph{chhāṛā}, once more
  \item \emph{Say it:} say \emph{chhāṛā}
  \item \emph{Recall:} say the Bengali for to lose, then the Bengali for to wear, then say \emph{chhāṛā} again
\end{itemize}

\begin{tcolorbox}[breakable,colback=teal!4,colframe=teal!35!black,title={Before you move on}]
What does \bn{ছাড়া} mean? (\textquotedblleft{}To leave, to let go\textquotedblright{}.) Say it once more.
\end{tcolorbox}
\section[\texorpdfstring{\bn{ফেলা} (phælā)}{ফেলা (phælā)}]{\bn{ফেলা} (phælā) — to throw, to drop}

\label{lesson:BN-C53-phla}



\begin{quote}
Before the new one: say the Bengali for to wear, then the Bengali for to leave.
\end{quote}

\subsection*{You'll want to know: \bn{ফেলা}}
\textbf{\bn{ফেলা}} — \emph{phælā} — \textquotedblleft{}to throw, to drop\textquotedblright{}.

\begin{cousinweb}
CDIAL lists \textbf{\bn{ফেলা}} (\emph{phælā}) under Sanskrit \emph{prathita}, \textquotedblleft{}spread out\textquotedblright{}, as a form extended with \emph{-ll-}, and asks whether it was crossed with \emph{*prēlayati}, \textquotedblleft{}pushes on\textquotedblright{}, a form it reconstructs (CDIAL 8651).
\end{cousinweb}

\begin{grammarlens}[title={What you've built}]
One more word for this chapter.
\end{grammarlens}

\subsection*{Your turn}
\begin{itemize}
\raggedright
  \item \emph{Say it:} \emph{phælā}
  \item \emph{Say it:} \emph{phælā}, once more
  \item \emph{Say it:} say \emph{phælā}
  \item \emph{Recall:} say the Bengali for to wear, then the Bengali for to leave, then say \emph{phælā} again
\end{itemize}

\begin{tcolorbox}[breakable,colback=teal!4,colframe=teal!35!black,title={Before you move on}]
What does \bn{ফেলা} mean? (\textquotedblleft{}To throw, to drop\textquotedblright{}.) Say it once more.
\end{tcolorbox}
\section[\texorpdfstring{\bn{ওড়া} (oṛā)}{ওড়া (oṛā)}]{\bn{ওড়া} (oṛā) — to fly}

\label{lesson:BN-C53-ora}



\begin{quote}
Before the new one: say the Bengali for to leave, then the Bengali for to throw.
\end{quote}

\subsection*{You'll want to know: \bn{ওড়া}}
\textbf{\bn{ওড়া}} — \emph{oṛā} — \textquotedblleft{}to fly\textquotedblright{}.

\begin{cousinweb}
\textbf{\bn{ওড়া}} (\emph{oṛā}) is a tadbhava, a word inherited from Sanskrit and reshaped by sound change. It comes from \emph{uḍḍīyate}, \textquotedblleft{}flies up\textquotedblright{} (CDIAL 1697).
\end{cousinweb}

\begin{grammarlens}[title={What you've built}]
One more word for this chapter.
\end{grammarlens}

\subsection*{Your turn}
\begin{itemize}
\raggedright
  \item \emph{Say it:} \emph{oṛā}
  \item \emph{Say it:} \emph{oṛā}, once more
  \item \emph{Say it:} say \emph{oṛā}
  \item \emph{Recall:} say the Bengali for to leave, then the Bengali for to throw, then say \emph{oṛā} again
\end{itemize}

\begin{tcolorbox}[breakable,colback=teal!4,colframe=teal!35!black,title={Before you move on}]
What does \bn{ওড়া} mean? (\textquotedblleft{}To fly\textquotedblright{}.) Say it once more.
\end{tcolorbox}
\section[\texorpdfstring{\bn{শেখানো} (shekhāno)}{শেখানো (shekhāno)}]{\bn{শেখানো} (shekhāno) — to teach}

\label{lesson:BN-C53-shekhano}



\begin{quote}
Before the new one: say the Bengali for to throw, then the Bengali for to fly.
\end{quote}

\subsection*{You'll want to know: \bn{শেখানো}}
\textbf{\bn{শেখানো}} — \emph{shekhāno} — \textquotedblleft{}to teach\textquotedblright{}.

\begin{cousinweb}
\textbf{\bn{শেখানো}} (\emph{shekhāno}) is a tadbhava, a word inherited from Sanskrit and reshaped by sound change. It comes from \emph{śikṣate}, \textquotedblleft{}learns\textquotedblright{} (CDIAL 12430). \emph{Shekhāno} is the causing form of \emph{shekhā}, \textquotedblleft{}to learn\textquotedblright{}; CDIAL lists the pair as \emph{sikhā} and \emph{sikhāna}.
\end{cousinweb}

\begin{grammarlens}[title={What you've built}]
One more word for this chapter.
\end{grammarlens}

\subsection*{Your turn}
\begin{itemize}
\raggedright
  \item \emph{Say it:} \emph{shekhāno}
  \item \emph{Say it:} \emph{shekhāno}, once more
  \item \emph{Say it:} say \emph{shekhāno}
  \item \emph{Recall:} say the Bengali for to throw, then the Bengali for to fly, then say \emph{shekhāno} again
\end{itemize}

\begin{tcolorbox}[breakable,colback=teal!4,colframe=teal!35!black,title={Before you move on}]
What does \bn{শেখানো} mean? (\textquotedblleft{}To teach\textquotedblright{}.) Say it once more.
\end{tcolorbox}
